% !Mode:: "TeX:UTF-8"

\usepackage{xeCJK}  
\usepackage{fontspec}
\usepackage{caption}
\usepackage{graphicx}
\usepackage{geometry}
%a4paper,text={150true mm,224true mm,head=5true mm,headsep=2.5true mm,foot=8.5true mm},
\geometry{top=3.8cm,bottom=3.8cm,left=3.2cm,right=3.2cm}

\usepackage{subfigure}
\usepackage{titlesec}               % 控制标题的宏包
\usepackage{titletoc}                   % 控制目录的宏包
\usepackage{fancyhdr}
\pagestyle{fancy}
\fancyhead{3.0cm}
\fancyfoot{3.0cm}                % fancyhdr宏包 页眉和页脚的相关定义
\usepackage[UTF8]{ctex}
\usepackage{color}          % 支持彩色
\usepackage{amsmath}        % AMSLaTeX宏包 用来排出更加漂亮的公式
\usepackage[below]{placeins}%允许上一个section的浮动图形出现在下一个section的开始部分,还提供\FloatBarrier命令,使所有未处理的浮动图形立即被处理
\usepackage{flafter}       % 使得所有浮动体不能被放置在其浮动环境之前，以免浮动体在引述它的文本之前出现.
\usepackage{multirow}       %使用Multirow宏包，使得表格可以合并多个row格
\usepackage{booktabs}       % 表格，横的粗线；\specialrule{1pt}{0pt}{0pt}
\usepackage{longtable}      %支持跨页的表格。
\usepackage{tabularx}
\usepackage{subfigure}%支持子图 %centerlast 设置最后一行是否居中
\usepackage[subfigure]{ccaption} %支持双语标题
\usepackage[sort&compress,numbers]{natbib}% 支持引用缩写的宏包
\usepackage{setspace}
\usepackage{enumitem}       %使用enumitem宏包,改变列表项的格式
\usepackage{calc}           %长度可以用+ - * / 进行计算
\usepackage{txfonts}
\usepackage{dsfont,amssymb,mathtools,amsfonts,mathrsfs,float}
\usepackage{bm}              % 处理数学公式中的黑斜体的宏包
\usepackage[amsmath,thmmarks,hyperref]{ntheorem}% 定理类环境宏包，其中 amsmath 选项用来兼容 AMS LaTeX 的宏包
\usepackage{natbib}
\setcitestyle{authoryear,round,colon}




\usepackage{appendix}
%\usepackage[toc,titletoc,title]{appendix}
\renewcommand{\appendixtocname}{附录}
\usepackage{chngcntr}

\usepackage{array}
\renewcommand\arraystretch{1.5}

\usepackage{amstext}
\usepackage{amsopn}
\usepackage{slashbox}
\pagestyle{empty}
\usepackage{pdfpages}
\usepackage{epsfig}
\usepackage{float}
\usepackage[singlelinecheck=off]{caption}
\usepackage[justification=centering]{caption}
%\captionsetup{justification=centering}
%\usepackage{subcaption}
%\captionsetup[table]{font=small,skip=0pt}
%\captionsetup[figure]{font=small,skip=0pt}
%\captionsetup{font={scriptsize}}
\usepackage{footnote}
\usepackage{threeparttable}

% 生成有书签的pdf及其开关, 该宏包应放在所有宏包的最后, 宏包之间有冲突
\def\atemp{dvipdfmx}\ifx\atemp\usewhat
\usepackage[dvipdfmx,unicode,           %dvi-->pdf 生成书签
            bookmarksnumbered=true,
            bookmarksopen=true,
            colorlinks=true,
            pdfborder={0 0 1},
            citecolor=black,
            linkcolor=black,
            anchorcolor=green,
            urlcolor=blue,
            breaklinks=true
            ]{hyperref}
\fi

\def\atemp{pdflatex}\ifx\atemp\usewhat
\usepackage{cmap}                       %pdflatex编译时，可以生成可复制、粘贴的中文PDF文档
\usepackage[pdftex,unicode,
            %CJKbookmarks=true,
            bookmarksnumbered=true,
            bookmarksopen=true,
            colorlinks=true,
            pdfborder={0 0 1},
            citecolor=black,
            linkcolor=black,
            anchorcolor=green,
            urlcolor=blue,
            breaklinks=true
            ]{hyperref}
\fi

\def\atempxetex{xelatex}\ifx\atempxetex\usewhat %\def\atempxetex{xelatex} main.tex中已定义;
\usepackage[xetex,
            bookmarksnumbered=true,
            bookmarksopen=true,
            colorlinks=true,
            pdfborder={0 0 1},
            citecolor=black,
            linkcolor=black,
            anchorcolor=green,
            urlcolor=blue,
            breaklinks=true,
            naturalnames  %与algorithm2e宏包协调
            ]{hyperref}
\fi

\usepackage[boxed,linesnumbered,algochapter]{algorithm2e}  % 算法的宏包，注意宏包兼容性，先后顺序为float、hyperref、algorithm(2e)，否则无法生成算法列表
